
\documentclass[10pt]{article}
\usepackage[utf8]{inputenc}
\usepackage[T1]{fontenc}
\usepackage{amsmath, amssymb, xcolor, geometry, enumitem, multicol, tcolorbox}

\definecolor{mainblue}{RGB}{0, 102, 204}
\geometry{a4paper, margin=1.2cm, top=1cm, bottom=3cm, footskip=1.2cm}

\begin{document}
\enlargethispage{2.5cm}

% Modern Header
\begin{tcolorbox}[colback=mainblue!5, colframe=mainblue, sharp corners, boxrule=0.5pt]
    \small \textbf{Unit:} Unit 0: Foundations of Algebra \hfill \textbf{Name:} \rule{3cm}{0.4pt} \\
    \small \textbf{Topic:} Inequalities \hfill \textbf{Date:} \rule{2cm}{0.4pt} \\
    \centering \textbf{\large Homework (Carolina Reaper)}
\end{tcolorbox}

\vspace{0.2cm}

% MCQ Section
\begin{multicols}{2}
\begin{enumerate}[label=\textbf{Q\arabic*.}, leftmargin=*, itemsep=6pt, parsep=0pt]

  \item \small Solve for $x$ in $2x - 5 > 11$.
  \begin{enumerate}[label=(\Alph*), leftmargin=0.6cm, nosep, itemsep=2pt]
    \item $x > 8$
    \item $x < 8$
    \item $x > 4$
    \item $x < 4$
    \item $x = 8$
  \end{enumerate}

  \item \small Graph the solution to $x + 2 geq 7$ on a number line.
  \begin{enumerate}[label=(\Alph*), leftmargin=0.6cm, nosep, itemsep=2pt]
    \item $[5, infty)$
    \item $(-infty, 5]$
    \item $[7, infty)$
    \item $(-infty, 7]$
    \item $[2, 7]$
  \end{enumerate}

  \item \small A music producer needs to mix a track with a frequency range of $200 leq f leq 800$ Hz. Write the solution in interval notation.
  \begin{enumerate}[label=(\Alph*), leftmargin=0.6cm, nosep, itemsep=2pt]
    \item $[200, 800]$
    \item $(200, 800)$
    \item $[200, 800)$
    \item $(200, 800]$
    \item $(-infty, 200] cup [800, infty)$
  \end{enumerate}

  \item \small An artist is creating a canvas with a width $w$ that satisfies $w - 3 > 10$. Solve for $w$.
  \begin{enumerate}[label=(\Alph*), leftmargin=0.6cm, nosep, itemsep=2pt]
    \item $w > 13$
    \item $w < 13$
    \item $w > 7$
    \item $w < 7$
    \item $w = 13$
  \end{enumerate}

  \item \small A rhythm pattern has a tempo $t$ that must satisfy $t + 2 leq 120$. Write the solution in interval notation.
  \begin{enumerate}[label=(\Alph*), leftmargin=0.6cm, nosep, itemsep=2pt]
    \item $(-infty, 118]$
    \item $[118, infty)$
    \item $[0, 118]$
    \item $[118, infty)$
    \item $(-infty, 118)$
  \end{enumerate}

  \item \small Solve for $x$ in $x/4 - 2 < 5$.
  \begin{enumerate}[label=(\Alph*), leftmargin=0.6cm, nosep, itemsep=2pt]
    \item $x < 28$
    \item $x > 28$
    \item $x < 14$
    \item $x > 14$
    \item $x = 28$
  \end{enumerate}

  \item \small A DJ needs to adjust the bass level $b$ to satisfy $b + 5 geq 20$. Write the solution in interval notation.
  \begin{enumerate}[label=(\Alph*), leftmargin=0.6cm, nosep, itemsep=2pt]
    \item $[15, infty)$
    \item $(-infty, 15]$
    \item $[15, 20]$
    \item $[20, infty)$
    \item $(-infty, 20]$
  \end{enumerate}

  \item \small Graph the solution to $3x - 2 leq 14$ on a number line.
  \begin{enumerate}[label=(\Alph*), leftmargin=0.6cm, nosep, itemsep=2pt]
    \item $(-infty, 16/3]$
    \item $[16/3, infty)$
    \item $[14/3, infty)$
    \item $(-infty, 14/3]$
    \item $[2, 14/3]$
  \end{enumerate}

\end{enumerate}
\end{multicols}

\vspace{-0.2cm}
\hrule
\vspace{0.2cm}
\textbf{\small Open Ended Questions (FRQ)}
\begin{enumerate}[label=\textbf{Q\arabic*.}, start=9, itemsep=5pt, topsep=0pt]

  \item \small A painter is creating a mural with a height $h$ that satisfies $h - 4 > 2h - 10$. Solve for $h$ and write the solution in interval notation. Show all work. \vspace{2cm}

  \item \small A sound engineer needs to adjust the treble level $t$ to satisfy $2t + 5 geq 11$. Write the solution in interval notation and graph the solution on a number line. Show all work. \vspace{2cm}

\end{enumerate}

\vfill

% Chef's Tips Footer
\begin{tcolorbox}[colback=blue!5, colframe=mainblue!50, title=\small \textbf{Chef's Tips}, fonttitle=\bfseries, bottom=1mm]
\begin{itemize}[noitemsep, topsep=2pt, leftmargin=0.5cm]
  \item \scriptsize Tip 1: Ensure students use the correct inequality signs when solving and graphing.\item \scriptsize Tip 2: Encourage students to check their work by plugging in values from the solution set.\item \scriptsize Tip 3: Remind students to label their graphs with the correct interval notation.
\end{itemize}
\end{tcolorbox}

\end{document}
  